\documentclass[11pt,a4paper]{article}

% ── Paquetes ──────────────────────────────────────────────────────────────
\usepackage[utf8]{inputenc}
\usepackage[T1]{fontenc}
\usepackage[spanish,es-tabla]{babel}
\usepackage{amsmath,amssymb,amsthm}
\usepackage{mathtools}
\usepackage{geometry}
\usepackage{enumitem}
\usepackage{booktabs}
\usepackage{array}
\usepackage{tikz}
\usetikzlibrary{automata,positioning,arrows.meta,babel}
\usepackage{hyperref}
\usepackage{xcolor}
\usepackage{framed}
\usepackage{tcolorbox}
\tcbuselibrary{theorems,breakable}

\geometry{margin=2.5cm}

% ── Entornos ──────────────────────────────────────────────────────────────
\theoremstyle{definition}
\newtheorem{definicion}{Definición}[section]
\newtheorem{teorema}{Teorema}[section]
\newtheorem{proposicion}{Proposición}[section]
\newtheorem{corolario}{Corolario}[section]
\newtheorem{lema}{Lema}[section]

\theoremstyle{remark}
\newtheorem{observacion}{Observación}[section]

\newtcolorbox{ejemplo}[1][]{%
  colback=blue!5!white,
  colframe=blue!60!black,
  fonttitle=\bfseries,
  title=Ejemplo: #1,
  breakable
}

\newtcolorbox{ejercicio}[1][]{%
  colback=green!5!white,
  colframe=green!50!black,
  fonttitle=\bfseries,
  title=Ejercicio propuesto: #1,
  breakable
}

\newtcolorbox{cajanota}[1][]{%
  colback=yellow!8!white,
  colframe=orange!70!black,
  fonttitle=\bfseries,
  title=Nota importante: #1,
  breakable
}

% ── Comandos útiles ───────────────────────────────────────────────────────
\newcommand{\R}{\mathbb{R}}
\newcommand{\N}{\mathbb{N}}
\newcommand{\E}{\mathbb{E}}
\newcommand{\Var}{\mathrm{Var}}
\newcommand{\Prob}{\mathbb{P}}
\newcommand{\e}{\mathrm{e}}
\DeclareMathOperator{\diag}{diag}

% ══════════════════════════════════════════════════════════════════════════
\title{%
  \textbf{Análisis Transitorio de Cadenas de Markov\\
  de Tiempo Continuo}\\[6pt]
  \large Lectura para curso de pregrado
}
\author{Notas de clase}
\date{2026}

\begin{document}
\maketitle

\tableofcontents
\newpage

% ======================================================================
%  SECCIÓN 1 — MOTIVACIÓN
% ======================================================================
\section{Motivación: ¿por qué nos interesa el comportamiento transitorio?}

Cuando estudiamos cadenas de Markov de tiempo continuo (CTMC), el primer
resultado que se enseña suele ser la \emph{distribución estacionaria}: el
vector $\boldsymbol{\pi}$ tal que $\boldsymbol{\pi}Q = \mathbf{0}$.  Ese
resultado describe el sistema ``a largo plazo'', pero ignora preguntas
prácticas que surgen constantemente:

\begin{itemize}[nosep]
  \item Si un servidor de atención al público abre a las 8:00 a.m.\ sin
        clientes, ¿cuál es la probabilidad de que a las 8:15 haya al menos
        dos personas en cola?
  \item Un enlace de red se acaba de recuperar de una falla.  ¿Cuánto
        tiempo se espera hasta que la carga supere su capacidad?
  \item Un hospital tiene dos camas de emergencia.  Ambas están libres a
        las 6:00 a.m.  ¿Cuál es la \emph{varianza} del número de camas
        ocupadas a las 7:00 a.m.?
\end{itemize}

Todas estas preguntas requieren conocer la distribución del proceso en un
instante \emph{particular} $t$, no solo cuando $t\to\infty$.  Eso es
precisamente el \textbf{análisis transitorio}.

% ======================================================================
%  SECCIÓN 2 — REPASO RÁPIDO DE CTMC
% ======================================================================
\section{Repaso rápido: cadena de Markov de tiempo continuo}

\begin{definicion}[CTMC]
Sea $\{X(t) : t \ge 0\}$ un proceso estocástico con espacio de estados
finito $\mathcal{S} = \{0,1,\dots,n-1\}$.  Decimos que es una
\textbf{cadena de Markov de tiempo continuo} si cumple la propiedad de
Markov:
\[
  \Prob\bigl(X(t+s) = j \mid X(s) = i,\; X(u),\; 0 \le u < s\bigr)
  = \Prob\bigl(X(t+s) = j \mid X(s) = i\bigr),
\]
y además esta probabilidad depende solo de la diferencia $t$ (proceso
homogéneo en el tiempo).
\end{definicion}

El comportamiento de la cadena queda determinado por su \textbf{matriz
generadora} (o \emph{rate matrix})
\[
  Q = \bigl[q_{ij}\bigr]_{i,j \in \mathcal{S}},
\]
donde $q_{ij} \ge 0$ para $i \ne j$ representa la tasa de transición de
$i$ a $j$, y la diagonal satisface
\[
  q_{ii} = -\sum_{j \ne i} q_{ij}.
\]
Esto garantiza que cada fila suma cero: $Q\mathbf{1} = \mathbf{0}$.

% ======================================================================
%  SECCIÓN 3 — ECUACIONES DE KOLMOGOROV
% ======================================================================
\section{Ecuaciones de Kolmogorov y la matriz de transición}

\subsection{La probabilidad de transición}

Definimos la \textbf{probabilidad de transición} como
\[
  p_{ij}(t) \;=\; \Prob\bigl(X(t) = j \mid X(0) = i\bigr),
  \qquad t \ge 0.
\]
Agrupando todas las entradas en una matriz $P(t) = [p_{ij}(t)]$, se
puede demostrar que $P(t)$ satisface las \emph{ecuaciones diferenciales
de Kolmogorov}.

\subsection{Ecuación de Kolmogorov hacia adelante (\textit{forward})}

\begin{teorema}[Ecuación de Kolmogorov \emph{forward}]
\[
  \frac{dP(t)}{dt} = P(t)\,Q, \qquad P(0) = I.
\]
\end{teorema}

\begin{proof}[Idea de la demostración]
Partimos de la propiedad de Chapman--Kolmogorov:
$P(t+h) = P(t)\,P(h)$.  Restando $P(t)$ a ambos lados:
\[
  P(t+h) - P(t) = P(t)\bigl[P(h) - I\bigr].
\]
Dividiendo por $h$ y tomando $h \to 0^+$, el lado derecho converge a
$P(t)\,Q$ porque $\lim_{h\to 0^+} \frac{P(h)-I}{h} = Q$ (definición del
generador).
\end{proof}

\subsection{Ecuación de Kolmogorov hacia atrás (\textit{backward})}

De forma análoga, escribiendo $P(t+h) = P(h)\,P(t)$ se obtiene
\[
  \frac{dP(t)}{dt} = Q\,P(t), \qquad P(0) = I.
\]

\begin{cajanota}[Ambas ecuaciones llevan a la misma solución]
Para espacio de estados finito y $Q$ acotada, ambas ecuaciones tienen la
misma solución única.  No siempre ocurre lo mismo en espacio infinito.
\end{cajanota}

% ======================================================================
%  SECCIÓN 4 — LA EXPONENCIAL DE MATRIZ
% ======================================================================
\section{La solución: la exponencial de matriz \texorpdfstring{$\e^{Qt}$}{exp(Qt)}}

\subsection{Definición y existencia}

La solución de $P'(t) = P(t)Q$ con $P(0)=I$ es, formalmente,
\[
  \boxed{P(t) = \e^{Qt} \;=\; \sum_{k=0}^{\infty} \frac{(Qt)^k}{k!}.}
\]

Esta serie converge absolutamente para toda matriz $Q$ y todo $t \ge 0$,
puesto que $\|Q\|$ es finita (espacio de estados finito).  Verifiquemos
que cumple la ecuación:
\[
  \frac{d}{dt}\e^{Qt}
  = \frac{d}{dt}\sum_{k=0}^{\infty}\frac{Q^k\,t^k}{k!}
  = \sum_{k=1}^{\infty}\frac{Q^k\,t^{k-1}}{(k-1)!}
  = Q\sum_{m=0}^{\infty}\frac{Q^m\,t^m}{m!}
  = Q\,\e^{Qt}.
\]

\subsection{Propiedades de \texorpdfstring{$\e^{Qt}$}{exp(Qt)}}

\begin{enumerate}[nosep]
  \item $\e^{Q \cdot 0} = I$.
  \item $\e^{Q(t+s)} = \e^{Qt}\,\e^{Qs}$ \quad (propiedad de semigrupo).
  \item Cada entrada de $\e^{Qt}$ está en $[0,1]$ y cada fila suma 1
        (es una matriz estocástica).
  \item $\lim_{t\to\infty}\e^{Qt}$ existe y sus filas convergen a
        $\boldsymbol{\pi}$ si la cadena es irreducible.
\end{enumerate}

% ──────────────────────────────────────────────────────────────────────
\subsection{Métodos para calcular \texorpdfstring{$\e^{Qt}$}{exp(Qt)}}
\label{sec:metodos}

A continuación presentamos varios métodos para calcular la exponencial
de la matriz generadora.  Cada uno tiene ventajas según el contexto.

% ── Método 1 ──────────────────────────────────────────────────────────
\subsubsection{Método 1: diagonalización}

Si $Q$ es \emph{diagonalizable}, existen una matriz invertible $V$ y una
matriz diagonal $\Lambda = \diag(\lambda_0,\lambda_1,\dots,\lambda_{n-1})$
tales que
\[
  Q = V \Lambda V^{-1}.
\]
Entonces
\[
  Q^k = V\Lambda^k V^{-1},
\]
y sustituyendo en la serie:
\[
  \e^{Qt} = V\,\e^{\Lambda t}\,V^{-1}
  = V \diag\!\bigl(\e^{\lambda_0 t},\e^{\lambda_1 t},\dots,
    \e^{\lambda_{n-1}t}\bigr)\,V^{-1}.
\]

\begin{cajanota}[Valores propios de $Q$]
Puesto que las filas de $Q$ suman cero, uno de los valores propios es
siempre $\lambda = 0$.  El resto tienen parte real estrictamente negativa
(si la cadena es irreducible), lo que asegura que
$\e^{\lambda_k t}\to 0$ cuando $t\to\infty$ para $k\ne 0$, recuperando
la distribución estacionaria.
\end{cajanota}

% ── Método 2 ──────────────────────────────────────────────────────────
\subsubsection{Método 2: serie truncada (definición directa)}

Para matrices pequeñas o para cálculos a mano con pocos términos:
\[
  \e^{Qt} \approx I + Qt + \frac{(Qt)^2}{2!} + \frac{(Qt)^3}{3!} + \cdots
\]
Se calcula potencia por potencia hasta que los términos son despreciables.

% ── Método 3 ──────────────────────────────────────────────────────────
\subsubsection{Método 3: uniformización (método de Jensen)}
\label{sec:uniformizacion}

Este método es el más usado en la práctica computacional.  La idea es
\emph{convertir} el problema de tiempo continuo en uno de tiempo discreto
con un reloj de Poisson.

Sea $\gamma \ge \max_i |q_{ii}|$ (la cota uniforme de las tasas de
salida).  Definimos la \textbf{matriz de transición discreta}
\[
  \tilde{P} = I + \frac{Q}{\gamma}.
\]
Se verifica fácilmente que $\tilde{P}$ es una matriz estocástica (entradas
no negativas, filas suman 1).

\begin{teorema}[Uniformización]
\[
  \e^{Qt} = \sum_{k=0}^{\infty} \e^{-\gamma t}\,
    \frac{(\gamma t)^k}{k!}\;\tilde{P}^k.
\]
\end{teorema}

\begin{proof}
Escribimos $Q = \gamma(\tilde{P}-I)$, de modo que
\begin{align*}
  \e^{Qt}
  &= \e^{\gamma(\tilde{P}-I)t}
   = \e^{-\gamma t}\;\e^{\gamma \tilde{P}\,t} \\
  &= \e^{-\gamma t}\sum_{k=0}^{\infty}\frac{(\gamma t)^k}{k!}\tilde{P}^k.
\end{align*}
\end{proof}

La interpretación probabilística es directa: $\e^{-\gamma t}\frac{(\gamma t)^k}{k!}$ es la probabilidad de que un proceso de Poisson de tasa $\gamma$ tenga exactamente $k$ eventos en $[0,t]$, y $\tilde{P}^k$ da las probabilidades de transición tras $k$ pasos de la cadena embebida discreta.

% ── Método 4 ──────────────────────────────────────────────────────────
\subsubsection{Método 4: transformada de Laplace}

Tomando la transformada de Laplace de $P'(t) = QP(t)$:
\[
  s\hat{P}(s) - I = Q\hat{P}(s)
  \quad\Longrightarrow\quad
  \hat{P}(s) = (sI - Q)^{-1}.
\]
Luego $P(t) = \mathcal{L}^{-1}\{(sI-Q)^{-1}\}$.  Útil cuando la
inversión se puede hacer analíticamente (matrices $2\times 2$).

% ── Método 5 ──────────────────────────────────────────────────────────
\subsubsection{Método 5: Cayley--Hamilton}

El teorema de Cayley--Hamilton dice que toda matriz satisface su propia
ecuación característica.  Para una matriz $n\times n$, la exponencial se
puede escribir como un polinomio de grado a lo más $n-1$ en $Q$:
\[
  \e^{Qt} = \alpha_0(t)\,I + \alpha_1(t)\,Q + \cdots
             + \alpha_{n-1}(t)\,Q^{n-1},
\]
donde los coeficientes $\alpha_k(t)$ se determinan evaluando en los
valores propios de $Q$.

Para una matriz $2\times 2$ con valores propios $\lambda_1, \lambda_2$
($\lambda_1 \ne \lambda_2$):
\[
  \e^{Qt} = \frac{\lambda_2\e^{\lambda_1 t}-\lambda_1\e^{\lambda_2 t}}
                  {\lambda_2-\lambda_1}\,I
           +\frac{\e^{\lambda_2 t}-\e^{\lambda_1 t}}
                  {\lambda_2-\lambda_1}\,Q.
\]

% ======================================================================
%  SECCIÓN 5 — DISTRIBUCIÓN EN UN INSTANTE t
% ======================================================================
\section{La distribución del estado en un instante particular \texorpdfstring{$t$}{t}}

\subsection{Vector de distribución transitoria}

Sea $\boldsymbol{\alpha} = (\alpha_0,\alpha_1,\dots,\alpha_{n-1})$ la
distribución inicial: $\alpha_i = \Prob(X(0)=i)$.  Entonces la
\textbf{distribución en el instante $t$} es
\[
  \boldsymbol{p}(t) = \boldsymbol{\alpha}\,\e^{Qt},
  \qquad p_j(t) = \Prob(X(t) = j).
\]

Si el sistema arranca en un estado determinista $X(0)=i$, entonces
$\boldsymbol{\alpha} = \mathbf{e}_i$ (el $i$-ésimo vector canónico) y
\[
  p_j(t) = \bigl[\e^{Qt}\bigr]_{ij}.
\]

\subsection{¿Qué preguntas podemos responder?}

Con $\boldsymbol{p}(t)$ a la mano, podemos calcular:

\begin{enumerate}
  \item \textbf{Probabilidad puntual:} $\Prob(X(t)=j)$ directamente.
  \item \textbf{Probabilidad de un conjunto:}
    $\Prob(X(t)\in A) = \sum_{j\in A} p_j(t)$.
  \item \textbf{Valor esperado:} si $f$ es una función sobre los estados,
    \[
      \E[f(X(t))] = \sum_{j\in\mathcal{S}} f(j)\,p_j(t).
    \]
  \item \textbf{Varianza:}
    \[
      \Var[f(X(t))] = \E[f(X(t))^2] - \bigl(\E[f(X(t))]\bigr)^2
      = \sum_j f(j)^2\,p_j(t) - \Bigl(\sum_j f(j)\,p_j(t)\Bigr)^{\!2}.
    \]
\end{enumerate}

\begin{cajanota}[Deducción de la fórmula de la varianza]
Recordemos que para una variable aleatoria discreta $Y$ con valores
$y_1,y_2,\dots$ y probabilidades $p_1,p_2,\dots$:
\begin{align*}
  \E[Y]   &= \sum_j y_j\,p_j, \\
  \E[Y^2] &= \sum_j y_j^2\,p_j, \\
  \Var[Y]  &= \E[Y^2] - (\E[Y])^2
            = \sum_j y_j^2\,p_j - \Bigl(\sum_j y_j\,p_j\Bigr)^{\!2}.
\end{align*}
En nuestro caso, $Y = f(X(t))$, $y_j = f(j)$ y $p_j = p_j(t)$.  No hay
nada ``nuevo'': es la fórmula estándar de la varianza aplicada a la
distribución transitoria.
\end{cajanota}

% ======================================================================
%  SECCIÓN 6 — EJEMPLO 1: DOS ESTADOS
% ======================================================================
\section{Ejemplo 1 --- Sistema con dos estados (proceso de nacimiento y muerte simple)}

\begin{ejemplo}[Un enlace de red que falla y se repara]
Un enlace de comunicación puede estar en uno de dos estados: funcionando
(estado~0) o caído (estado~1).  Cuando funciona, las fallas ocurren a
una tasa de $\lambda = 2$ fallas/hora.  Cuando está caído, las
reparaciones se completan a una tasa de $\mu = 3$ reparaciones/hora.  El
enlace acaba de ser puesto en servicio (arranca en estado~0) a las 9:00~a.m.

Queremos saber: ¿cuál es la probabilidad de que esté caído a las
9:20~a.m.?  ¿Cuál es el número esperado de ``enlaces caídos'' (en un
sentido trivial: 0 o 1) a esa hora y su varianza?

\tcblower

\textbf{Paso 1: definir $Q$.}

El diagrama de transición es:

\begin{center}
\begin{tikzpicture}[->,>=Stealth,node distance=3cm,
    every state/.style={thick,fill=blue!10}]
  \node[state] (0) {0};
  \node[state] (1) [right=of 0] {1};
  \path (0) edge[bend left] node[above] {$\lambda=2$} (1)
        (1) edge[bend left] node[below] {$\mu=3$} (0);
\end{tikzpicture}
\end{center}

La matriz generadora es
\[
  Q = \begin{pmatrix} -2 & 2 \\ 3 & -3 \end{pmatrix}.
\]

\textbf{Paso 2: calcular $\e^{Qt}$ por diagonalización.}

El polinomio característico es
\[
  \det(Q-\lambda I) = (-2-\lambda)(-3-\lambda) - 6
  = \lambda^2 + 5\lambda = \lambda(\lambda+5).
\]
Los valores propios son $\lambda_1 = 0$ y $\lambda_2 = -5$.

Para $\lambda_1=0$: $(Q-0\cdot I)\mathbf{v}=\mathbf{0}$ da
$\mathbf{v}_1 = \begin{pmatrix}1\\1\end{pmatrix}$ (núcleo de $Q$).

Para $\lambda_2=-5$: $(Q+5I)\mathbf{v}=\mathbf{0}$ da
$\begin{pmatrix}3&2\\3&2\end{pmatrix}\mathbf{v}=\mathbf{0}$, así que
$\mathbf{v}_2 = \begin{pmatrix}2\\-3\end{pmatrix}$.

Formamos $V = \begin{pmatrix}1&2\\1&-3\end{pmatrix}$, con
$V^{-1} = \frac{1}{-5}\begin{pmatrix}-3&-2\\-1&1\end{pmatrix}
= \begin{pmatrix}3/5&2/5\\1/5&-1/5\end{pmatrix}$.

Entonces:
\begin{align*}
  \e^{Qt} &= V\begin{pmatrix}1&0\\0&\e^{-5t}\end{pmatrix}V^{-1} \\
  &= \begin{pmatrix}1&2\\1&-3\end{pmatrix}
     \begin{pmatrix}1&0\\0&\e^{-5t}\end{pmatrix}
     \begin{pmatrix}3/5&2/5\\1/5&-1/5\end{pmatrix} \\
  &= \frac{1}{5}\begin{pmatrix}
       3+2\e^{-5t} & 2-2\e^{-5t} \\
       3-3\e^{-5t} & 2+3\e^{-5t}
     \end{pmatrix}.
\end{align*}

\textbf{Verificaciones rápidas:}
\begin{itemize}[nosep]
  \item $t=0$: $\e^{Q\cdot 0}=\frac{1}{5}\begin{pmatrix}5&0\\0&5\end{pmatrix}=I$. \checkmark
  \item $t\to\infty$: $\e^{Qt}\to\frac{1}{5}\begin{pmatrix}3&2\\3&2\end{pmatrix}$
    $\Rightarrow$ $\boldsymbol{\pi}=(3/5,\;2/5)$. \checkmark
  \item Cada fila suma 1 para todo $t$. \checkmark
\end{itemize}

\textbf{Paso 3: evaluar en $t = 1/3$ hora (20 minutos).}

Con $\boldsymbol{\alpha}=(1,0)$ (arranca en estado~0):
\[
  \boldsymbol{p}(1/3) = (1,0)\,\e^{Q/3}
  = \Bigl(\frac{3+2\e^{-5/3}}{5},\;\frac{2-2\e^{-5/3}}{5}\Bigr).
\]
Numéricamente, $\e^{-5/3} \approx 0.18888$, así que
\[
  p_0(1/3) \approx \frac{3+0.3776}{5} = \frac{3.3776}{5} \approx 0.6755,
  \qquad
  p_1(1/3) \approx \frac{2-0.3776}{5} = \frac{1.6224}{5} \approx 0.3245.
\]

La probabilidad de que el enlace esté caído a las 9:20 a.m.\ es
aproximadamente $32.45\%$.

\textbf{Paso 4: valor esperado y varianza del número de enlaces caídos.}

Definimos $f(j) = j$ (el estado mismo indica cuántos enlaces están caídos:
0 o 1).

\[
  \E[X(t)] = 0\cdot p_0(t) + 1\cdot p_1(t) = p_1(t)
  = \frac{2-2\e^{-5t}}{5}.
\]
En $t=1/3$:
\[
  \E[X(1/3)] \approx 0.3245.
\]

Para la varianza:
\[
  \E[X(t)^2] = 0^2\cdot p_0(t) + 1^2\cdot p_1(t) = p_1(t),
\]
ya que $f(j)^2 = j^2$ y los estados son 0 y 1.  Por lo tanto
\begin{align*}
  \Var[X(t)] &= \E[X(t)^2] - (\E[X(t)])^2 \\
  &= p_1(t) - p_1(t)^2 \\
  &= p_1(t)\bigl(1-p_1(t)\bigr).
\end{align*}
Numéricamente:
\[
  \Var[X(1/3)] \approx 0.3245 \times 0.6755 \approx 0.2192.
\]

Observe que la varianza tiene la misma forma que la de una Bernoulli, lo
cual es natural: $X(t)$ toma solo los valores $\{0,1\}$, así que
$X(t)\sim\mathrm{Bernoulli}(p_1(t))$.
\end{ejemplo}

% ======================================================================
%  SECCIÓN 7 — EJEMPLO 2: COLA M/M/1/2
% ======================================================================
\section{Ejemplo 2 --- Sistema de cola con capacidad limitada}

\begin{ejemplo}[Ventanilla de trámites con sala de espera pequeña]
Una oficina municipal tiene una sola ventanilla para trámites.  El área
de espera solo admite a una persona adicional (aparte de la que se está
atendiendo), así que la capacidad total del sistema es de 2 personas.  Los
ciudadanos llegan con tasa $\lambda = 4$ personas/hora, y el empleado
atiende con tasa $\mu = 5$ personas/hora.  Si un ciudadano llega y el
sistema está lleno (2 personas), se va sin ser atendido.

La oficina abre a las 8:00 a.m.\ vacía.  Nos preguntan:

\begin{enumerate}[label=(\alph*)]
  \item ¿Cuál es la distribución del número de personas en el sistema a
        las 8:30 a.m.?
  \item ¿Cuál es el número esperado de personas y su varianza a esa hora?
\end{enumerate}

\tcblower

\textbf{Paso 1: modelo y generadora.}

Este es un sistema $M/M/1/2$ (notación de Kendall).  El espacio de estados
es $\mathcal{S}=\{0,1,2\}$ y el diagrama de transiciones es:

\begin{center}
\begin{tikzpicture}[->,>=Stealth,node distance=2.5cm,
    every state/.style={thick,fill=blue!10}]
  \node[state] (0) {0};
  \node[state] (1) [right=of 0] {1};
  \node[state] (2) [right=of 1] {2};
  \path (0) edge[bend left] node[above] {$4$} (1)
        (1) edge[bend left] node[below] {$5$} (0)
        (1) edge[bend left] node[above] {$4$} (2)
        (2) edge[bend left] node[below] {$5$} (1);
\end{tikzpicture}
\end{center}

\[
  Q = \begin{pmatrix}
    -4 & 4 & 0 \\
    5 & -9 & 4 \\
    0 & 5 & -5
  \end{pmatrix}.
\]

\textbf{Paso 2: $\e^{Qt}$ por uniformización.}

Para una matriz $3\times 3$ la diagonalización manual es posible pero
tediosa.  Usemos uniformización (Sección~\ref{sec:uniformizacion}).

Elegimos $\gamma = 9$ (el máximo de $|q_{ii}|$).  La matriz discreta es
\[
  \tilde{P} = I + \frac{Q}{9}
  = \begin{pmatrix}
    5/9 & 4/9 & 0 \\
    5/9 & 0   & 4/9 \\
    0   & 5/9 & 4/9
  \end{pmatrix}.
\]

La fórmula de uniformización da
\[
  \boldsymbol{p}(t) = \boldsymbol{\alpha}\,\e^{Qt}
  = \sum_{k=0}^{\infty}\e^{-9t}\frac{(9t)^k}{k!}
    \;\boldsymbol{\alpha}\,\tilde{P}^k.
\]

Con $\boldsymbol{\alpha}=(1,0,0)$ y $t=0.5$ (media hora), tenemos
$\gamma t = 4.5$.  Calculemos los primeros términos manualmente para
ver cómo funciona.

\textbf{$k=0$:}
$\boldsymbol{\alpha}\tilde{P}^0 = (1,0,0)$.

\textbf{$k=1$:}
$\boldsymbol{\alpha}\tilde{P} = (5/9,\;4/9,\;0)$.

\textbf{$k=2$:}
\[
  \boldsymbol{\alpha}\tilde{P}^2
  = (5/9,\;4/9,\;0)\tilde{P}
  = \Bigl(\frac{25+20}{81},\;\frac{20+0}{81},\;\frac{16}{81}\Bigr)
  = \Bigl(\frac{45}{81},\;\frac{20}{81},\;\frac{16}{81}\Bigr).
\]

Los pesos de Poisson son $w_k = \e^{-4.5}\frac{4.5^k}{k!}$:
\begin{align*}
  w_0 &= \e^{-4.5} \approx 0.01111, \\
  w_1 &= 4.5\,\e^{-4.5} \approx 0.04999, \\
  w_2 &= \frac{4.5^2}{2}\,\e^{-4.5} \approx 0.11248, \\
  w_3 &\approx 0.16872, \\
  w_4 &\approx 0.18981, \quad\text{etc.}
\end{align*}

En la práctica, se suman términos hasta que la cola de Poisson acumulada
restante es menor que una tolerancia (p.~ej., $10^{-8}$).  Para este
ejemplo, calculemos el resultado numéricamente:

\[
  \boldsymbol{p}(0.5) \approx (0.5007,\;0.3203,\;0.1790).
\]

(Estos valores se obtienen de la evaluación numérica completa de la serie
o de la exponencial de matriz directamente.)

\textbf{Paso 3: valor esperado.}

Sea $N = X(0.5)$ el número de personas en el sistema.  Definimos
$f(j)=j$.
\[
  \E[N] = 0 \times 0.5007 + 1 \times 0.3203 + 2 \times 0.1790
  = 0.3203 + 0.3580 = 0.6783.
\]

\textbf{Paso 4: varianza.}

Primero, el segundo momento:
\[
  \E[N^2] = 0^2\times 0.5007 + 1^2\times 0.3203 + 2^2\times 0.1790
  = 0.3203 + 0.7160 = 1.0363.
\]
Luego,
\[
  \Var[N] = 1.0363 - (0.6783)^2 = 1.0363 - 0.4601 = 0.5762.
\]

\textbf{Paso 5: interpretación.}

A las 8:30 a.m., hay un $50\%$ de probabilidad de que la ventanilla esté
vacía, un $32\%$ de que haya exactamente una persona, y un $18\%$ de que
el sistema esté lleno.  En promedio hay $0.68$ personas, con desviación
estándar $\sqrt{0.5762}\approx 0.759$.

Comparando con el estado estable: la distribución estacionaria
de un $M/M/1/2$ con $\rho=\lambda/\mu=4/5$ es
\[
  \pi_j = \frac{1-\rho}{1-\rho^3}\rho^j,
  \qquad j=0,1,2.
\]
Numéricamente $\boldsymbol{\pi}\approx(0.4098,\;0.3279,\;0.2623)$, con
$\E_\pi[N]\approx 0.8525$.  Así que a los 30 minutos el sistema aún no
ha llegado al estado estable: hay más probabilidad de estar vacío y
menos probabilidad de estar lleno que en régimen.
\end{ejemplo}

% ======================================================================
%  SECCIÓN 8 — EJEMPLO CON DOS VARIABLES
% ======================================================================
\section{Ejemplo 3 --- Análisis transitorio con dos variables}

\begin{ejemplo}[Dos servidores independientes]
Una compañía de soporte técnico tiene dos líneas telefónicas
independientes.  La línea~A recibe llamadas con tasa $\lambda_A=1$/min y
las resuelve con tasa $\mu_A=2$/min.  La línea~B recibe llamadas con tasa
$\lambda_B=2$/min y las resuelve con tasa $\mu_B=3$/min.  Cada línea
puede atender a lo más una llamada (las llamadas extra se pierden).

Las dos líneas se modelan como CTMC independientes.

A las 10:00 a.m., ambas líneas están libres.  Se pide:

\begin{enumerate}[label=(\alph*)]
  \item Distribución conjunta del par $(X_A(t),X_B(t))$ en $t=0.5$ min.
  \item Valor esperado y varianza del número total de líneas ocupadas
        $N(t) = X_A(t)+X_B(t)$ en $t=0.5$ min.
\end{enumerate}

\tcblower

\textbf{Paso 1: análisis de cada línea por separado.}

Cada línea es un sistema de dos estados ($\{0,1\}$), igual al
Ejemplo~1.

\medskip
\emph{Línea A:}
\[
  Q_A = \begin{pmatrix}-1&1\\2&-2\end{pmatrix},
  \quad
  \lambda_1=0,\;\lambda_2=-3.
\]
Usando la fórmula del caso $2\times 2$ (Sección~\ref{sec:metodos},
Cayley--Hamilton):
\[
  \e^{Q_A t} = \frac{1}{3}\begin{pmatrix}
    2+\e^{-3t} & 1-\e^{-3t} \\
    2-2\e^{-3t} & 1+2\e^{-3t}
  \end{pmatrix}.
\]
En $t=0.5$: $\e^{-1.5}\approx 0.2231$.
\[
  \boldsymbol{p}_A(0.5) = (1,0)\e^{Q_A/2}
  = \Bigl(\frac{2+0.2231}{3},\;\frac{1-0.2231}{3}\Bigr)
  \approx (0.7410,\;0.2590).
\]

\emph{Línea B:}
\[
  Q_B = \begin{pmatrix}-2&2\\3&-3\end{pmatrix},
  \quad
  \lambda_1=0,\;\lambda_2=-5.
\]
\[
  \e^{Q_B t} = \frac{1}{5}\begin{pmatrix}
    3+2\e^{-5t} & 2-2\e^{-5t} \\
    3-3\e^{-5t} & 2+3\e^{-5t}
  \end{pmatrix}.
\]
En $t=0.5$: $\e^{-2.5}\approx 0.08209$.
\[
  \boldsymbol{p}_B(0.5) = (1,0)\e^{Q_B/2}
  \approx \Bigl(\frac{3+0.1642}{5},\;\frac{2-0.1642}{5}\Bigr)
  \approx (0.6328,\;0.3672).
\]

\textbf{Paso 2: distribución conjunta.}

Dado que las líneas son \textbf{independientes}, la distribución conjunta
es el producto:
\[
  \Prob(X_A=a,\;X_B=b;\;t) = p_A^{(a)}(t)\cdot p_B^{(b)}(t).
\]

\begin{center}
\begin{tabular}{c|cc|c}
\toprule
$X_A\backslash X_B$ & 0 & 1 & Total \\
\midrule
0 & $0.7410\times 0.6328 = 0.4689$ & $0.7410\times 0.3672 = 0.2721$ & $0.7410$ \\
1 & $0.2590\times 0.6328 = 0.1639$ & $0.2590\times 0.3672 = 0.0951$ & $0.2590$ \\
\midrule
Total & $0.6328$ & $0.3672$ & $1$ \\
\bottomrule
\end{tabular}
\end{center}

\textbf{Paso 3: distribución de $N = X_A + X_B$.}

$N$ toma valores en $\{0,1,2\}$:
\begin{align*}
  \Prob(N=0) &= \Prob(X_A=0,X_B=0) = 0.4689, \\
  \Prob(N=1) &= \Prob(X_A=1,X_B=0) + \Prob(X_A=0,X_B=1) = 0.1639+0.2721 = 0.4360, \\
  \Prob(N=2) &= \Prob(X_A=1,X_B=1) = 0.0951.
\end{align*}

\textbf{Paso 4: valor esperado de $N$.}

\[
  \E[N] = 0\times 0.4689 + 1\times 0.4360 + 2\times 0.0951
  = 0.4360 + 0.1902 = 0.6262.
\]

Equivalentemente, por la linealidad de la esperanza:
\[
  \E[N] = \E[X_A] + \E[X_B] = 0.2590 + 0.3672 = 0.6262.
\]

\textbf{Paso 5: varianza de $N$.}

Puesto que $X_A$ y $X_B$ son \emph{independientes}:
\[
  \Var[N] = \Var[X_A] + \Var[X_B].
\]
Cada variable es Bernoulli, así que $\Var[X_i] = p_i(1-p_i)$:
\begin{align*}
  \Var[X_A] &= 0.2590\times 0.7410 = 0.1919, \\
  \Var[X_B] &= 0.3672\times 0.6328 = 0.2323, \\
  \Var[N]   &= 0.1919 + 0.2323 = 0.4242.
\end{align*}

\textbf{Paso 6: cálculo integral equivalente.}

Queremos verificar $\E[N]$ usando la expresión integral que surge de
la formulación ``desde primeros principios''.  En general,
\[
  \E[N(t)] = \sum_{j=0}^{2} j\cdot\Prob(N(t)=j).
\]
Pero supongamos que quisiéramos calcular $\E[X_A(t)]$ directamente a
partir de la solución de Kolmogorov.  Sabemos que
$p_1^A(t) = \frac{1-\e^{-3t}}{3}$, de modo que
\[
  \E[X_A(t)] = p_1^A(t) = \frac{1-\e^{-3t}}{3}.
\]
Si evaluamos numéricamente la integral de la ecuación diferencial
$\frac{dp_1^A}{dt} = q_{01}\,p_0^A(t) + q_{11}\,p_1^A(t)
= 1\cdot(1-p_1^A)-2\,p_1^A = 1 - 3p_1^A$ con $p_1^A(0)=0$, la
solución es exactamente $\frac{1-\e^{-3t}}{3}$ y en $t=0.5$ obtenemos
$0.2590$, consistente con lo anterior.

Alternativamente, podemos interpretar $\E[X_A(t)]$ como una integral de
la tasa de acumulación:
\[
  \E[X_A(t)] = \int_0^{t}\frac{d}{ds}\E[X_A(s)]\,ds
  = \int_0^{t}\bigl(1 - 3\,\E[X_A(s)]\bigr)\,ds.
\]
Resolviendo la ODE $m'(t) = 1-3m(t)$, $m(0)=0$:
\[
  m(t) = \frac{1}{3}(1-\e^{-3t}).
\]
\textbf{Numéricamente, integrando con un paso de Euler con $\Delta t=0.1$:}
\begin{center}
\begin{tabular}{ccc}
\toprule
$t$ & $m(t)$ Euler & $m(t)$ exacto \\
\midrule
$0.0$ & $0.0000$ & $0.0000$ \\
$0.1$ & $0.1000$ & $0.0950$ \\
$0.2$ & $0.1700$ & $0.1647$ \\
$0.3$ & $0.2190$ & $0.2148$ \\
$0.4$ & $0.2533$ & $0.2498$ \\
$0.5$ & $0.2773$ & $0.2590$ \\
\bottomrule
\end{tabular}
\end{center}
(El método de Euler con $\Delta t = 0.1$ sobreestima ligeramente; con
pasos más finos converge al valor exacto.)
\end{ejemplo}

% ======================================================================
%  SECCIÓN 9 — EJEMPLO 4: COLA M/M/1/2 CON INTEGRALES DETALLADAS
% ======================================================================
\section{Ejemplo 4 --- Cálculos integrales detallados en una cola $M/M/1/1$}

\begin{ejemplo}[Consultorio médico con un solo paciente]
Un consultorio médico solo admite a un paciente a la vez (no hay sala de
espera).  Los pacientes llegan con tasa $\lambda=3$/hora y el médico
atiende con tasa $\mu=6$/hora.  Si alguien llega y el consultorio está
ocupado, se va.  El consultorio abre vacío.

Calculemos $\E[X(t)]$ y $\Var[X(t)]$ para $t$ arbitrario, y evaluemos
numéricamente en $t=10$ minutos $= 1/6$ hora.

\tcblower

\textbf{La generadora:}
\[
  Q = \begin{pmatrix}-3&3\\6&-6\end{pmatrix}.
\]
Valores propios: $0$ y $-9$.

\textbf{Exponencial de matriz:}
\[
  \e^{Qt} = \frac{1}{9}\begin{pmatrix}
    6+3\e^{-9t} & 3-3\e^{-9t} \\
    6-6\e^{-9t} & 3+6\e^{-9t}
  \end{pmatrix}.
\]
Con $\boldsymbol{\alpha}=(1,0)$:
\[
  p_1(t) = \frac{3-3\e^{-9t}}{9} = \frac{1-\e^{-9t}}{3}.
\]

\textbf{Valor esperado:}
\[
  \E[X(t)] = p_1(t) = \frac{1-\e^{-9t}}{3}.
\]

\textbf{Varianza:}
Como $X(t)\in\{0,1\}$:
\[
  \Var[X(t)] = p_1(t)(1-p_1(t))
  = \frac{1-\e^{-9t}}{3}\cdot\frac{2+\e^{-9t}}{3}
  = \frac{(1-\e^{-9t})(2+\e^{-9t})}{9}.
\]

\textbf{Deducción por integración directa.}

La ecuación diferencial para $p_1(t)$ (Kolmogorov forward) es:
\[
  p_1'(t) = 3\,p_0(t) - 6\,p_1(t) = 3(1-p_1(t)) - 6\,p_1(t)
  = 3 - 9\,p_1(t).
\]
Esta es una ODE lineal de primer orden.  Resolviéndola por factor
integrante:

Multiplicamos por $\e^{9t}$:
\[
  \e^{9t}p_1'(t) + 9\e^{9t}p_1(t) = 3\e^{9t}
  \quad\Longleftrightarrow\quad
  \frac{d}{dt}\bigl[\e^{9t}p_1(t)\bigr] = 3\e^{9t}.
\]
Integramos de $0$ a $t$:
\[
  \e^{9t}p_1(t) - \underbrace{p_1(0)}_{=0}
  = 3\int_0^t \e^{9s}\,ds
  = 3\cdot\frac{\e^{9t}-1}{9}
  = \frac{\e^{9t}-1}{3}.
\]
\[
  \boxed{p_1(t) = \frac{1-\e^{-9t}}{3}.}
\]

\textbf{Evaluación numérica de la integral.}

Comprobemos numéricamente que $\int_0^{1/6}\e^{9s}\,ds$ da el valor
correcto:
\[
  \int_0^{1/6}\e^{9s}\,ds = \frac{\e^{9/6}-1}{9} = \frac{\e^{3/2}-1}{9}
  \approx \frac{4.4817-1}{9} = \frac{3.4817}{9} \approx 0.3869.
\]
Entonces
\[
  p_1(1/6) = \e^{-3/2}\times\frac{\e^{3/2}-1}{3}
  = \frac{1-\e^{-3/2}}{3}
  \approx \frac{1-0.2231}{3} = \frac{0.7769}{3} \approx 0.2590.
\]

\[
  \E[X(1/6)] \approx 0.2590, \qquad
  \Var[X(1/6)] \approx 0.2590\times 0.7410 \approx 0.1919.
\]

\textbf{Aproximación numérica con regla del trapecio.}

Para visualizar el cálculo integral, usemos la regla del trapecio con
$n=3$ subintervalos en $[0,1/6]$ (paso $h=1/18$):
\begin{align*}
  s_0 &= 0,\quad s_1 = 1/18,\quad s_2 = 2/18,\quad s_3 = 3/18 = 1/6.\\
  \e^{9s_0} &= 1, \quad
  \e^{9/18} = \e^{0.5} \approx 1.6487, \quad
  \e^{1} \approx 2.7183, \quad
  \e^{1.5} \approx 4.4817.
\end{align*}
Regla del trapecio:
\[
  \int_0^{1/6}\e^{9s}\,ds \approx \frac{h}{2}
  \bigl[f(s_0) + 2f(s_1) + 2f(s_2) + f(s_3)\bigr]
  = \frac{1/18}{2}[1+3.2974+5.4366+4.4817]
  \approx 0.3954.
\]
Comparando con el valor exacto $0.3869$: el error relativo es
$\sim 2.2\%$, típico para solo 3 subintervalos.  Con más subintervalos
converge rápidamente.
\end{ejemplo}

% ======================================================================
%  SECCIÓN 10 — EJEMPLO CON DOS VARIABLES Y COLAS
% ======================================================================
\section{Ejemplo 5 --- Dos colas acopladas y la distribución conjunta}

\begin{ejemplo}[Dos cajeros en un banco con desbordamiento]
Un banco pequeño tiene dos cajeros.  El cajero~1 es rápido ($\mu_1=6$/h)
y el cajero~2 es estándar ($\mu_2=4$/h).  Los clientes llegan al
banco con tasa total $\lambda=5$/h.  Un cliente que llega va al cajero~1
si está libre; si no, va al cajero~2 si está libre; si ambos están
ocupados, se va.

El estado del sistema es el par $(C_1,C_2)$ donde $C_i\in\{0,1\}$
indica si el cajero $i$ está ocupado.  El espacio de estados es
$\mathcal{S}=\{(0,0),(1,0),(0,1),(1,1)\}$, que enumeramos como estados
$0,1,2,3$ respectivamente.

El banco abre a las 9:00 a.m.\ con ambos cajeros libres.

\tcblower

\textbf{Paso 1: construir $Q$.}

Las transiciones son:
\begin{itemize}[nosep]
  \item $(0,0)\to(1,0)$: llega un cliente, va al cajero~1. Tasa $\lambda=5$.
  \item $(1,0)\to(0,0)$: el cajero~1 termina. Tasa $\mu_1=6$.
  \item $(1,0)\to(1,1)$: llega un cliente, cajero~1 ocupado, va al~2. Tasa $\lambda=5$.
  \item $(0,1)\to(1,1)$: llega un cliente, cajero~2 ocupado, va al~1. Tasa $\lambda=5$.
  \item $(0,1)\to(0,0)$: el cajero~2 termina. Tasa $\mu_2=4$.
  \item $(1,1)\to(0,1)$: el cajero~1 termina. Tasa $\mu_1=6$.
  \item $(1,1)\to(1,0)$: el cajero~2 termina. Tasa $\mu_2=4$.
\end{itemize}

\[
  Q = \begin{pmatrix}
    -5 & 5  & 0 & 0  \\
    6  & -11 & 0 & 5  \\
    4  & 0  & -9 & 5  \\
    0  & 4  & 6  & -10
  \end{pmatrix}.
\]

\textbf{Paso 2: calcular $\e^{Qt}$ numéricamente.}

Para una matriz $4\times 4$, es más práctico usar el método de
uniformización.  Tomemos $\gamma = 11$ y formemos
$\tilde{P} = I + Q/11$.

En lugar de detallar todos los pasos (que son mecánicos), presentamos
el resultado numérico para $t=1/6$ hora (10 minutos), usando
suficientes términos de la serie de uniformización:

Con $\boldsymbol{\alpha}=(1,0,0,0)$:
\[
  \boldsymbol{p}(1/6) \approx (0.5122,\;0.3185,\;0.0570,\;0.1123).
\]

\textbf{Paso 3: preguntas sobre la variable $N = C_1+C_2$ (total de
cajeros ocupados).}

$N$ toma valores $\{0,1,2\}$:
\begin{align*}
  \Prob(N=0) &= p_{(0,0)} = 0.5122, \\
  \Prob(N=1) &= p_{(1,0)} + p_{(0,1)} = 0.3185 + 0.0570 = 0.3755, \\
  \Prob(N=2) &= p_{(1,1)} = 0.1123.
\end{align*}

\[
  \E[N] = 0\times 0.5122 + 1\times 0.3755 + 2\times 0.1123 = 0.6001.
\]

\[
  \E[N^2] = 0 + 0.3755 + 4\times 0.1123 = 0.8247.
\]

\[
  \Var[N] = 0.8247 - (0.6001)^2 = 0.8247 - 0.3601 = 0.4646.
\]

\textbf{Paso 4: preguntas por cajero.}

También podemos calcular las marginales:
\begin{align*}
  \Prob(C_1=1) &= p_{(1,0)} + p_{(1,1)} = 0.3185 + 0.1123 = 0.4308, \\
  \Prob(C_2=1) &= p_{(0,1)} + p_{(1,1)} = 0.0570 + 0.1123 = 0.1693.
\end{align*}

Note que $\E[N] = \E[C_1]+\E[C_2] = 0.4308+0.1693 = 0.6001$. \checkmark

Para la varianza necesitamos la covarianza (porque $C_1$ y $C_2$ \textbf{no
son independientes} en este modelo):
\begin{align*}
  \E[C_1 C_2] &= \Prob(C_1=1,C_2=1) = 0.1123, \\
  \mathrm{Cov}(C_1,C_2) &= \E[C_1 C_2]-\E[C_1]\E[C_2]
  = 0.1123 - 0.4308\times 0.1693 \\
  &= 0.1123 - 0.0730 = 0.0393.
\end{align*}
\[
  \Var[N] = \Var[C_1]+\Var[C_2]+2\,\mathrm{Cov}(C_1,C_2).
\]
$\Var[C_1] = 0.4308\times 0.5692 = 0.2452$, \quad
$\Var[C_2] = 0.1693\times 0.8307 = 0.1406$.
\[
  \Var[N] = 0.2452+0.1406+2(0.0393) = 0.4644 \approx 0.4646.
\]
(La pequeña diferencia es redondeo.)  Observe cómo la covarianza positiva
\emph{aumenta} la varianza de $N$ respecto al caso independiente.
\end{ejemplo}



% ======================================================================
%  BIBLIOGRAFÍA SUGERIDA
% ======================================================================
\section*{Bibliografía sugerida}

\begin{enumerate}[nosep]
  \item W.~J.~Stewart, \textit{Probability, Markov Chains, Queues, and
        Simulation}, Princeton University Press, 2009.
        Capítulos sobre CTMC y análisis transitorio.
  \item G.~Bolch, S.~Greiner, H.~de~Meer, K.~S.~Trivedi,
        \textit{Queueing Networks and Markov Chains}, 2nd ed., Wiley, 2006.
        Métodos numéricos para exponenciales de matriz.
  \item C.~Moler, C.~Van~Loan, ``Nineteen Dubious Ways to Compute the
        Exponential of a Matrix, Twenty-Five Years Later'',
        \textit{SIAM Review}, 45(1):3--49, 2003.
        Referencia clásica sobre métodos para $\e^{At}$.
  \item R.~Nelson, \textit{Probability, Stochastic Processes, and Queueing
        Theory}, Springer, 1995.
        Exposición accesible de CTMCs y colas.
\end{enumerate}

\end{document}
